%% Sections shared by every manual, English edition.
%% Pulled in with \input{common-en}. The only thing that differs between
%% applets is \appletfolder, which is also the folder name and the path on the site.

\section{Publishing it for a class}

Any static host will do: GitHub Pages, Netlify, a university web directory,
even a shared folder served over HTTP. The only thing to upload is the
\code{index.html} file, inside a folder named however you want the address to
read.

The repository ships a GitHub Actions workflow
(\code{.github/workflows/pages.yml}) that publishes every applet at once on each
push to the default branch. To switch it on, in the repository:
\emph{Settings\,$\rightarrow$\,Pages\,$\rightarrow$\,Source: GitHub Actions}.
That is the one step the workflow cannot do for itself, because creating a Pages
site needs administration rights that a \code{GITHUB\_TOKEN} is never granted.

\begin{amnote}
To hand the applet to students without reliable internet, send them the
\code{index.html} file directly. It weighs less than a photograph and opens with
a double-click.
\end{amnote}

\section{Licence and provenance}

The application code is MIT. Use it, change it and pass it on, in class or out
of it, with attribution.

This applet is a standalone rewrite of an original GeoGebra applet. The rewrite
is not a literal translation: where the original had a construction error, a
division by zero reachable with its own sliders, or a dead object inherited from
whatever applet it was copied from, this version fixes it and says so. The
section \emph{Changes from the GeoGebra original} lists those changes one by
one.

Everything is hand-written and dependency-free: the expression parser, the
symbolic differentiation, the level-curve tracing and the drawing. In
particular there is no \code{eval}, so what a student types into a text box
cannot become code, and a strict content security policy does not break the
page.
